\section{Related Work}

Much work has gone into designing environments to test and develop MARL agents. However, not many of these focused on providing a qualitatively
challenging environment that would provide together elements of partial
observability, challenging dynamics, and
high-dimensional observation spaces.

\citet{stone2005keepaway} presented Keepaway soccer, a domain built on the RoboCup soccer
simulator \citep{kitano1997robocup}, a 2D simulation of a football environment
with simplified physics, where the main task consists of keeping a ball within a
pre-defined area where agents in teams can reach, steal, and pass the ball,
providing a simplified setup for studying cooperative MARL.
This domain was later extended to the Half Field Offense task
\citep{kalyanakrishnan2006half, hausknecht_half_2016}, which increases the
difficulty of the problem by requiring the agents to not only keep the ball
within bounds but also to score a goal.
Neither task scales well in difficulty with the number of agents, as most
agents need to do little coordination. There is also a lack of interesting environment
dynamics beyond the simple 2D physics nor good reward signals, thus reducing the
impact of the environment as a testbed.

Multiple gridworld-like environments have also been explored.
\citet{lowe_multi-agent_2017} released a set of simple grid-world like
environments for multi-agent RL alongside an implementation of MADDPG, featuring
a mix of competitive and cooperative tasks focused on shared communication and
low level continuous control.
\citet{leibo_multi-agent_2017} show several mixed-cooperative Markov
environment focused on testing social dilemmas, however, they did not release an
implementation to further explore the tasks.
\citet{yang2018mean, zheng2017magent} present a framework for creating
gridworlds focuses on many-agents tasks, where the number of agents ranges from
the hundreds to the millions. This work, however, focuses on testing for emergent
behaviour, since environment dynamics and control space need to remain
relatively simple for the tasks to be tractable.
\citet{resnick2018pommerman} propose a multi-agent environment based on the game
\emph{Bomberman}, encompassing a series of cooperative and adversarial tasks
meant to provide a more challenging set of tasks with a relatively
straightforward 2D state observation and simple grid-world-like action spaces.

Learning to play StarCraft games also has been investigated in several
communities: work ranging from evolutionary algorithms to tabular
RL applied has shown that the game is an excellent testbed
for both modelling and planning \citep{ontanon2013survey}, however, most have
focused on single-agent settings with multiple controllers and classical
algorithms.
More recently, progress has been made on developing frameworks that enable
researchers working with deep neural networks to test recent algorithms on these
games; work on applying deep RL algorithms to single-agent and multi-agent
versions of the micromanagement tasks has thus been steadily appearing
\citep{usunier2016episodic, foerster_stabilising_2017,
foerster_counterfactual_2017, rashid_qmix:_2018, nardelli2018value,
hu2018knowledge, shao2018starcraft, foerster2018multi} with the release of
TorchCraft \citep{synnaeve_torchcraft:_2016} and SC2LE
\citep{vinyals_starcraft_2017}, interfaces to respectively \emph{StarCraft:
BroodWar} and \emph{StarCraft II}. Our work presents the first standardised
testbed for decentralised control in this space.

Other work focuses on playing the full game of StarCraft, including 
macromanagement and tactics 
\citep{pang2018reinforcement,sun2018tstarbots,alphastar}. 
By introducing decentralisation and local observability, our agents 
are excessively restricted compared to normal full gameplay. SMAC is therefore 
not intended as an environment to train agents for use in full gameplay. 
Instead, we use the StarCraft II game engine to build rich and interesting 
multi-agent problems.